\chapter{2019年北京卷（文）}

\section{单选题}

\begin{problem}
已知集合 \(A=\{x\mid -1<x<2\}\)，\(B=\{x\mid x>1\}\)，则 \(A\cup B=\)（\quad）.

\choices
    {\(\left(-1,1\right)\)}
    {\(\left(1,2\right)\)}
    {\(\left(-1,+\infty\right)\)}
    {\(\left(1,+\infty\right)\)}
\end{problem}

\begin{answer}
C
\end{answer}

\begin{problem}
已知复数 \(z=2+\i\)，则 \(z\overline z=\)（\quad）.

\choices
    {\(\sqrt3\)}
    {\(\sqrt5\)}
    {\(3\)}
    {\(5\)}
\end{problem}

\begin{answer}
D
\end{answer}

\begin{problem}
下列函数中，在区间 \(\left(0,+\infty\right)\) 上单调递增的是（\quad）.

\choices
    {\(y=\sqrt{x}\)}
    {\(y=2^{-x}\)}
    {\(y=\log_{\frac12}x\)}
    {\(y=\dfrac1x\)}
\end{problem}

\begin{answer}
A
\end{answer}

\begin{problem}
执行如图所示的程序框图，输出的 \(s\) 值为

\texfigure[max width=6.0cm]{img_repaint/2019/beijing_liberal/q4_fig1.tex}

\choices
    {\(1\)}
    {\(2\)}
    {\(3\)}
    {\(4\)}
\end{problem}

\begin{answer}
B
\end{answer}

\begin{problem}
已知双曲线 \(\dfrac{x^2}{a^2}-y^2=1\)（\(a>0\)）的离心率是 \(\sqrt5\)，则 \(a=\)（\quad）.

\choices
    {\(6\)}
    {\(4\)}
    {\(2\)}
    {\(\dfrac12\)}
\end{problem}

\begin{answer}
D
\end{answer}

\begin{problem}
设函数 \(f(x)=\cos x+b\sin x\)（\(b\) 为常数），则“\(b=0\)”是“\(f(x)\) 为偶函数”的（\quad）.

\choices
    {充分而不必要条件}
    {必要而不充分条件}
    {充分必要条件}
    {既不充分也不必要条件}
\end{problem}

\begin{answer}
C
\end{answer}

\begin{problem}
在天文学中，天体的明暗程度可以用星等或亮度来描述. 两颗星的星等与亮度满足 \(m_2-m_1=\dfrac52\lg\dfrac{E_1}{E_2}\)，其中星等为 \(m_k\) 的星的亮度为 \(E_k\)（\(k=1\)，\(2\)）. 已知太阳的星等是 \(-26.7\)，天狼星的星等是 \(-1.45\)，则太阳与天狼星的亮度的比值为（\quad）.

\choices
    {\(10^{10.1}\)}
    {\(10.1\)}
    {\(\lg 10.1\)}
    {\(10^{-10.1}\)}
\end{problem}

\begin{answer}
A
\end{answer}

\begin{problem}
如图，\(A\)，\(B\) 是半径为 \(2\) 的圆周上的定点，\(P\) 为圆周上的动点，\(\angle APB\) 是锐角，大小为 \(\beta\). 图中阴影区域的面积的最大值为

\texfigure[max width=6.0cm]{img_repaint/2019/beijing_liberal/q8_fig1.tex}

\choices
    {\(4\beta+4\cos\beta\)}
    {\(4\beta+4\sin\beta\)}
    {\(2\beta+2\cos\beta\)}
    {\(2\beta+2\sin\beta\)}
\end{problem}

\begin{answer}
B
\end{answer}

\section{填空题}

\begin{problem}
已知向量 \(\bs a=(-4,3)\)，\(\bs b=(6,m)\)，且 \(\bs a\perp\bs b\)，则 \(m=\fillinblank{}\).
\end{problem}

\begin{answer}
\(8\)
\end{answer}

\begin{problem}
若 \(x\)，\(y\) 满足 \(x\le2\)，\(y\ge-1\)，\(4x-3y+1\ge0\)，则 \(y-x\) 的最小值为\fillinblank{}，最大值为\fillinblank{}.
\end{problem}

\begin{answer}
\(-3\)，\(1\)
\end{answer}

\begin{problem}
设抛物线 \(y^2=4x\) 的焦点为 \(F\)，准线为 \(l\). 则以 \(F\) 为圆心，且与 \(l\) 相切的圆的方程为\fillinblank{}.
\end{problem}

\begin{answer}
\((x-1)^2+y^2=4\)
\end{answer}

\begin{problem}
某几何体是由一个正方体去掉一个四棱柱所得，其三视图如图所示. 如果网格纸上小正方形的边长为 \(1\)，那么该几何体的体积为\fillinblank{}.

\texfigure[max width=0.58\linewidth]{img_repaint/2019/beijing_liberal/q12_fig1.tex}
\end{problem}

\begin{answer}
\(40\)
\end{answer}

\begin{problem}
已知 \(l\)，\(m\) 是平面 \(\alpha\) 外的两条不同直线. 给出下列三个论断：

\begin{circlelist}
\item \(l\perp m\)；

\item \(m\parallel\alpha\)；

\item \(l\perp \alpha\).
\end{circlelist}

以其中的两个论断作为条件，余下的一个论断作为结论，写出一个正确的命题：\fillinblank{}.
\end{problem}

\begin{answer}
若 \(l\perp m\)，\(l\perp\alpha\)，则 \(m\parallel\alpha\).
\end{answer}

\begin{problem}
李明自主创业，在网上经营一家水果店，销售的水果中有草莓，京白梨，西瓜，桃，价格依次为 \(60\) 元/盒，\(65\) 元/盒，\(80\) 元/盒，\(90\) 元/盒. 为增加销量，李明对这四种水果进行促销：一次购买水果的总价达到 \(120\) 元，顾客就少付 \(x\) 元. 每笔订单顾客网上支付成功后，李明会得到支付款的 \(80\%\).

\begin{enumerate}
\item 当 \(x=10\) 时，顾客一次购买草莓和西瓜各 \(1\) 盒，需要支付\fillinblank{} 元；

\item 在促销活动中，为保证李明每笔订单得到的金额均不低于促销前总价的七折，则 \(x\) 的最大值为\fillinblank{}.
\end{enumerate}
\end{problem}

\begin{answer}
\(130\)，\(15\)
\end{answer}

\section{解答题}

\begin{problem}
在 \(\triangle ABC\) 中，\(a=3\)，\(b-c=2\)，\(\cos B=-\dfrac12\).

\begin{enumerate}
\item 求 \(b\)，\(c\) 的值；

\item 求 \(\sin(B+C)\) 的值.
\end{enumerate}
\end{problem}

\begin{solution}
\begin{enumerate}
\item \(b=7\)，\(c=5\).
\item \(\sin(B+C)=\dfrac{3\sqrt3}{14}\).
\end{enumerate}
\end{solution}

\begin{problem}
设 \(\{a_n\}\) 是等差数列，\(a_1=-10\)，且 \(a_2+10\)，\(a_3+8\)，\(a_4+6\) 成等比数列.

\begin{enumerate}
\item 求 \(\{a_n\}\) 的通项公式；

\item 记 \(\{a_n\}\) 的前 \(n\) 项和为 \(S_n\)，求 \(S_n\) 的最小值.
\end{enumerate}
\end{problem}

\begin{solution}
\begin{enumerate}
\item \(a_n=2n-12\).
\item \(S_n\) 的最小值为 \(S_6=-30\).
\end{enumerate}
\end{solution}

\begin{problem}
改革开放以来，人们的支付方式发生了巨大转变. 近年来，移动支付已成为主要支付方式之一. 为了解某校学生上个月 \(A\)，\(B\) 两种移动支付方式的使用情况，从全校所有的 \(1000\) 名学生中随机抽取了 \(100\) 人，发现样本中 \(A\)，\(B\) 两种支付方式都不使用的有 \(5\) 人，样本中仅使用 \(A\) 和仅使用 \(B\) 的学生的支付金额分布情况如下：

\begin{center}
\begin{tabular}{|c|c|c|}
\hline
支付方式 & 不大于 \(2\,000\) 元 & 大于 \(2\,000\) 元 \\
\hline
仅使用 \(A\) & \(27\) 人 & \(3\) 人 \\
\hline
仅使用 \(B\) & \(24\) 人 & \(1\) 人 \\
\hline
\end{tabular}
\end{center}

\begin{enumerate}
\item 估计该校学生中上个月 \(A\)，\(B\) 两种支付方式都使用的人数；

\item 从样本仅使用 \(B\) 的学生中随机抽取 \(1\) 人，求该学生上个月支付金额大于 \(2\,000\) 元的概率；

\item 已知上个月样本学生的支付方式在本月没有变化. 现从样本仅使用 \(B\) 的学生中随机抽查 \(1\) 人，发现他本月的支付金额大于 \(2\,000\) 元. 结合（2）的结果，能否认为样本仅使用 \(B\) 的学生中本月支付金额大于 \(2\,000\) 元的人数有变化？说明理由.
\end{enumerate}
\end{problem}

\begin{solution}
\begin{enumerate}
\item 估计该校学生中上个月 \(A\)，\(B\) 两种支付方式都使用的人数为 \(400\) 人.
\item 该学生上个月支付金额大于 \(2\,000\) 元的概率为 \(\dfrac1{25}\).
\item 答案不唯一：可以认为有变化，也可以认为无法确定是否变化.
\end{enumerate}
\end{solution}

\begin{problem}
如图，在四棱锥 \(P-ABCD\) 中，\(PA\perp\) 平面 \(ABCD\)，底面 \(ABCD\) 为菱形，\(E\) 为 \(CD\) 的中点.

\begin{enumerate}
\item 求证：\(BD\perp\) 平面 \(PAC\)；

\item 若 \(\angle ABC=60^\circ\)，求证：平面 \(PAB\perp\) 平面 \(PAE\)；

\item 棱 \(PB\) 上是否存在点 \(F\)，使得 \(CF\parallel\) 平面 \(PAE\)？说明理由.
\end{enumerate}

\texfigure[max width=0.52\linewidth]{img_repaint/2019/beijing_liberal/q18_fig1.tex}
\end{problem}

\begin{solution}
\begin{enumerate}
\item 证明 \(BD\perp\) 平面 \(PAC\).
\item 证明平面 \(PAB\perp\) 平面 \(PAE\).
\item 存在点 \(F\)，且 \(F\) 为棱 \(PB\) 的中点，使得 \(CF\parallel\) 平面 \(PAE\).
\end{enumerate}
\end{solution}

\begin{problem}
已知椭圆 \(C:\dfrac{x^2}{a^2}+\dfrac{y^2}{b^2}=1\) 的右焦点为 \((1,0)\)，且经过点 \(A(0,1)\).

\begin{enumerate}
\item 求椭圆 \(C\) 的方程；

\item 设 \(O\) 为原点，直线 \(l:y=kx+t\)（\(t\ne\pm1\)）与椭圆 \(C\) 交于两个不同点 \(P\)，\(Q\)，直线 \(AP\) 与 \(x\) 轴交于点 \(M\)，直线 \(AQ\) 与 \(x\) 轴交于点 \(N\)，若 \(\lvert OM\rvert\cdot\lvert ON\rvert=2\)，求证：直线 \(l\) 经过定点.
\end{enumerate}
\end{problem}

\begin{solution}
\begin{enumerate}
\item 椭圆 \(C\) 的方程为 \(\dfrac{x^2}{2}+y^2=1\).
\item 直线 \(l\) 经过定点 \((0,0)\).
\end{enumerate}
\end{solution}

\begin{problem}
已知函数 \(f(x)=\dfrac14x^3-x^2+x\).

\begin{enumerate}
\item 求曲线 \(y=f(x)\) 的斜率为 \(1\) 的切线方程；

\item 当 \(x\in\left[-2,4\right]\) 时，求证：\(x-6\le f(x)\le x\)；

\item 设 \(F(x)=\lvert f(x)-(x+a)\rvert\)（\(a\in\R\)），记 \(F(x)\) 在区间 \(\left[-2,4\right]\) 上的最大值为 \(M(a)\)，当 \(M(a)\) 最小时，求 \(a\) 的值.
\end{enumerate}
\end{problem}

\begin{solution}
\begin{enumerate}
\item 切线方程为 \(y=x\) 或 \(y=x-\dfrac{64}{27}\).
\item 当 \(x\in\left[-2,4\right]\) 时，\(x-6\le f(x)\le x\).
\item 当 \(M(a)\) 最小时，\(a=-3\).
\end{enumerate}
\end{solution}
